%% LyX 2.4.4 created this file.  For more info, see https://www.lyx.org/.
%% Do not edit unless you really know what you are doing.
\documentclass[english]{article}
\usepackage[T1]{fontenc}
\usepackage[latin9]{inputenc}
\usepackage{enumitem}
\usepackage{amsmath}
\usepackage{amsthm}
\usepackage{amssymb}
\usepackage{geometry}
\geometry{verbose,tmargin=1in,bmargin=1in,lmargin=1in,rmargin=1in}

\makeatletter
%%%%%%%%%%%%%%%%%%%%%%%%%%%%%% Textclass specific LaTeX commands.
\numberwithin{equation}{section}
\numberwithin{figure}{section}
\newlength{\lyxlabelwidth}      % auxiliary length 

%%%%%%%%%%%%%%%%%%%%%%%%%%%%%% User specified LaTeX commands.
\usepackage{fancyhdr}
\usepackage{lastpage}
\pagestyle{fancy}
\usepackage{enumitem}
\usepackage{ifthen}
%sets math fonts to be more readable
\everymath=\expandafter{\the\everymath\displaystyle}
%\DeclareMathSizes{10}{10}{10}{10}
\usepackage{multicol}
% Added by lyx2lyx
\setlength{\parskip}{\smallskipamount}
\setlength{\parindent}{0pt}

\makeatother

\usepackage{babel}
\begin{document}
\renewcommand{\author}{Dr.\,James Ashe}

\newcommand{\classNum}{131}

\newcommand{\classSec}{B and I}

\newcommand{\examType}{Midterm Review}

\renewcommand{\date}{October 30, 2012}

%%First page's header%%%%%%%%%%%%%%%%%%%%%%
\fancypagestyle{firststyle} %%header settings for first pagr
{
	\fancyhead[L]{MTH \classNum{}(\classSec{}) \author{}\\
	\textbf{\examType{}} \date{}}
	\fancyhead[C]{\textbf{}}
	\setlength{\headsep}{14pt}
}
%%%%%%%%%%%%%%%%%%%%%%%%%%%%%%%%%%%%%%%%%%

%%Next page's header%%%%%%%%%%%%%%%%%%%%%%%
\setlist{nolistsep}
\lhead{\classNum{}(\classSec{})} 
\chead{\textbf{}} 
\cfoot{\thepage{}~of~\pageref{LastPage}} 
\setlength{\headsep}{5pt} 
%%%%%%%%%%%%%%%%%%%%%%%%%%%%%%%%%%%%%%%%%%%

%%Various settings; see the preamble for other settings%%%%%
%%\setenumerate[0]{leftmargin=12pt,itemindent=0pt,labelwidth=0pt}
\renewcommand{\labelenumi}{\textbf{\arabic{enumi}.}}
\renewcommand{\labelenumii}{\textbf{\alph{enumii}.}}
\thispagestyle{firststyle}
%%%%%%%%%%%%%%%%%%%%%%%%%%%%%%%%%%%%%%%%%%

\textbf{Find the $x$ and $y$ intercepts from a graph. }(See the
midterm review for examples)
\begin{enumerate}
\item Find all the times the graph intersects the $x-$axis. These are your
$x-$intercepts.
\item Find all the times the graph intersects the $y-$axis. These are your
$y-$intercepts.
\end{enumerate}
\textbf{Solve an equation with fractions. ex. $\frac{2x}{6}+5=\frac{3x}{9}$}
\begin{enumerate}
\item Multiply by the Lowest Common Denominator (LCD). In this case, $18$.
Note BOTH terms on the left got multiplied by $18$.
\begin{eqnarray*}
18\left(\frac{2x}{6}+2\right) & = & \left(\frac{x}{9}\right)18\\
\frac{36x}{6}+36 & = & \frac{18x}{9}\\
6x+36 & = & 2x
\end{eqnarray*}
\item Now just solve.
\begin{eqnarray*}
4x & = & -36\\
x & = & -\frac{36}{4}=-9
\end{eqnarray*}
\end{enumerate}
\textbf{Solve a formula for a specified variable. ex. $P=2L+2W$ }for
$L$.
\begin{enumerate}
\item Just solve like you do for an $x$, and treat the other letters like
numbers. 
\begin{eqnarray*}
P & = & 2L+2W\\
P-2W & = & 2L\\
\frac{P-2W}{2} & = & L
\end{eqnarray*}
\item Note: The $2$ divides BOTH factors, and the you DO NOT cancel anything
in $2W$.
\end{enumerate}
\textbf{Solve using the square root property. ex. $(x-4)^{2}+1=2$.}
\begin{enumerate}
\item Get the squared part by itself:
\[
(x-4)^{2}=1
\]
\item Use a square root; Don't forget the $\pm$. 
\begin{eqnarray*}
\sqrt{(x-4)^{2}} & = & \pm\sqrt{1}\\
x-4 & = & \pm1
\end{eqnarray*}
\item Solve.
\begin{eqnarray*}
x & = & \pm1+4\\
x=5 & \mbox{ and } & x=3
\end{eqnarray*}
\end{enumerate}
\textbf{Solve a radical. ex. $\sqrt{x-1}+2=5$.}
\begin{enumerate}
\item Get the radical by itself, 
\begin{eqnarray*}
\sqrt{x-1}+2 & = & 5\\
\sqrt{x-1} & = & 3
\end{eqnarray*}
\item Use a square root, simplify, and solve. 
\begin{eqnarray*}
\left(\sqrt{x-1}\right)^{2} & = & 3^{2}\\
x-1 & = & 9\\
x & = & 10
\end{eqnarray*}
\item Check! $\sqrt{10-1}=9\mbox{\ensuremath{\checkmark}}$\newpage
\end{enumerate}
\textbf{Solve the absolute value or indicate that is has no solution.
ex. $|x+2|-2=10$}
\begin{enumerate}
\item Get the absolute value by itself.
\begin{eqnarray*}
\left|x+2\right|-2 & = & 10\\
\left|x+2\right| & = & 12
\end{eqnarray*}
\item Solve for $+$ and $-$.
\begin{eqnarray*}
x+2=-12 & \mbox{ and } & x+2=12\\
x=-14 &  & x=10
\end{eqnarray*}
\item Check! $\left|-14+2\right|-2=10\checkmark$ and $\left|10+2\right|-2=10\checkmark$
\end{enumerate}
\textbf{Solve by factoring. ex. $3x^{2}=-3x+6$}
\begin{enumerate}
\item Set the equation equal to $0$. 
\[
3x^{2}+x-2=0
\]
\item Factor any common multiples.
\[
3(x^{2}+x-2)=0
\]
\item Factor and solve.
\begin{eqnarray*}
3(x+2)(x-1) & = & 0\\
x=-2 & \mbox{ and } & x=1
\end{eqnarray*}
\end{enumerate}
\textbf{Solve the equation using the quadratic formula. ex. $3x^{2}=-2x+2$.}
\begin{enumerate}
\item Set the equation equal to $0$.
\[
3x^{2}+2x-2=0
\]
\item Identify the $a,b,c$ in $ax^{2}+bx+c=0$
\[
a=3,b=2,c=-2
\]
\item Use the quadratic formula, $x=\frac{-b\pm\sqrt{b^{2}-4ac}}{2a}$,
and simplify.
\[
x=\frac{-2\pm\sqrt{(2)^{2}-4(3)(-2)}}{2(3)}=\frac{-2\pm\sqrt{28}}{6}=\frac{-1\pm\sqrt{7}}{3}
\]
\end{enumerate}
\textbf{Solve the polynomial by factoring and using the zero principle.
ex. $x^{3}=-x^{2}+2x$.}
\begin{enumerate}
\item Set the equation equal to $0$.
\[
x^{3}+x^{2}-2x=0
\]
\item Factor out as many $x's$ as you can.
\[
x(x^{2}+x-2)=0
\]
\item Factor. 
\[
x(x+2)(x-1)=0
\]
\item Solve. DO NOT forget the $x=0$ part!
\[
x=0,x=-2,x=1
\]
\end{enumerate}
\textbf{Use the 5 step strategy for solving the word problem. When
3 times the number is subtracted 7 times the number, the result is
44.}
\begin{enumerate}
\item Summarize: ``3 times a number'' means, $3x$. ``7 times a number'',
means $7x$. ``Subtracted from'' means, $7x-3x$. ``Result is 44''
means, $=44$.
\item Write an equation: $7x-3x=44$
\item Solve: $x=11$.
\item Check: $7(11)-3(11)=44\checkmark$
\item State your answer as a sentence (this isn't applicable here).\newpage
\end{enumerate}
\textbf{Solve the word problem. ex. }A car rental place charges \$300
a week to rent a car plus an additional \$0.50 per mile. How many
miles can you drive in a week for \$400.
\begin{enumerate}
\item Write as an equation.
\[
300+0.50x=400
\]
\item Solve. 
\begin{eqnarray*}
0.50x & = & 100\\
x & = & \frac{100}{0.50}=\$200
\end{eqnarray*}
\end{enumerate}
\textbf{Find the Domain and the Range.}
\begin{enumerate}
\item Domain$=x's$; Range$=y's$.
\item From a table or list, just find the $x's$ and $y's$. ex. 
\begin{eqnarray*}
\left\{ \left(1,2\right),\left(0,3\right),\left(-1,4\right),\left(-2,5\right)\right\} \\
\mbox{domain=} & \left\{ 1,0,-1,-2\right\} \\
\mbox{range=} & \left\{ 2,3,4,5\right\} 
\end{eqnarray*}
\item From a graph, look at the horizontal values for the domain ($x's$)
and the vertical values for the range ($y's$). If they appear to
go of the graph, then use $\pm\infty$ (See the midterm review for
examples). 
\end{enumerate}
\textbf{Determine if the relation is a function. ex. $\left\{ \left(1,-2\right),\left(2,-1\right),\left(2,0\right),(3,1)\right\} $}
\begin{enumerate}
\item Check that each $x$ , in $\left(x,y\right)$, is paired with only
one $y$ value.
\begin{eqnarray*}
1 & \rightarrow & -2\\
2 & \rightarrow & -1\\
2 & \rightarrow & 0\\
3 & \rightarrow & 1
\end{eqnarray*}
\item If any $x$ is paired with two different $y's$ then the relation
is NOT a function.
\begin{eqnarray*}
2 & \rightarrow & -1\\
2 & \rightarrow & 0
\end{eqnarray*}
 So this is NOT a function.
\item Be suspicious if an $x$ appears more than once.
\end{enumerate}
\textbf{Find the slope passing through the given points. ex, $(1,3)$
}and $(0,5)$
\begin{enumerate}
\item Use the formula, $m=\frac{y_{2}-y_{1}}{x_{2}-x_{1}}$. Identify $(x_{1},y_{1})$
and $(x_{2},y_{2})$.
\[
x_{1}=1,\mbox{ }y_{1}=3\mbox{ and \ensuremath{x_{2}=0,\mbox{ }y_{2}=5}}
\]
\item Plug the values into the formula.
\[
m=\frac{5-3}{0-1}=\frac{2}{-1}=-2
\]
\end{enumerate}
\textbf{Use the information to write an equation for the line in point-slope
formula. ex. }Slope of $-2$, passing through $(1,3)$.
\begin{enumerate}
\item Use the formula, $y-y_{1}=m(x-x_{1})$. Identify $m$, $x_{1},$and
$y_{1}$. \\
\[
\mbox{The slope is \ensuremath{m=2}, and the point gives \ensuremath{x_{1}=1} and \ensuremath{y_{1}=3}}.
\]
\item Plug the values into the formula.
\[
y-3=-2(x-1)
\]
\item Simplify and write as $y=mx+b$.
\begin{eqnarray*}
y & = & -2x+2+3\\
y & = & -2x+5
\end{eqnarray*}
\end{enumerate}

\end{document}
